Ch 3 

\subsection{人工智能科学} 

人工智能(ArtificialIntelligence，AI)始于一个简单的想法。
计算机被发明之后，人们很快认识到它的能力远不限于数值计算，可以通过对信息和指令编程完成很多以往只有人脑才能完成的任务。


强调像人一样
强调合理性
思考

像人一样地思考

合理地思考

行动
像人一样地行动
合理地行动


||强调像人一样|强调合理性|
|:----:|:----:|:----:|
|思考|像人一样地思考|合理地思考|
|行动|像人一样地思考|合理地行动|


\subsection{人工智能：广跨各学科，综合文理工}

人工智能：对人的意识、思维过程进行模拟的一门学科。

\subsection{发展历程：AI的三次热潮和相应的研究学科突破}

人工智能发展如同过山车，忽上忽下，从创立以来不断分化

AI数理基础切换:从基于逻辑的表达与推理到基于统计的学习与计算

\subsection{三次热潮的总结}

第一阶段(前30年)：以数理逻辑的表达与推理为主。
杰出代表，如JohnMcCarthy、MarvinMinsky、HerbertSimmon：懂很多认知科学的东西，有很强的全局观念，拿过图灵奖和很多大奖。但工具基本都是基于数理逻辑和推理。
逻辑的东西发展得干净、漂亮。但符号的知识表达不落地；


 第二阶段(近30年)：以概率统计的建模、学习和计算为主。
视觉、自然语言、认知、机器学习、机器人五大领域在1990年中期都找到了概率统计这个新“机制”：统计建模、机器学习、随机算法……
